\section{Discussion and Limitations}
\label{sec:discussion}
\paragraph{Unresolvable evidence and finite hypotheses.}
Unknown scattering and missed candidate modes limit risk interpretation. More points or denoising cannot replace absent measurements. Preserve intervals or coherent alternatives when needed, and evaluate error at matched coverage.

\paragraph{Task and domain scope.}
The initial system supports a typed library of sensing quantities. Architecture reuse across indoor and automotive domains does not establish a universal pretrained model. Multi-person interactions, severe occlusion, and fast motion require explicit evaluation rather than inclusion by assumption.

\paragraph{Deployment and adaptation.}
Calibration may drift under new devices, layouts, or scattering conditions. Quantify required recalibration and training effort. If VLM control offers no benefit beyond simple interfaces for the supported tasks, narrow the language claim accordingly.

\paragraph{Data, privacy, and application boundaries.}
Radar data and reconstructed body geometry can still be sensitive. Follow dataset permissions and document any new participant collection and consent. Specify training-only cameras/depth sensors and the inference modality. These experiments do not establish clinical validity or driving safety.

\paragraph{Publication planning.}
This manuscript and the ICCV outline describe one system and share an evidence plan. Choose one submission based on the completed contribution and then-current official instructions. Keep full result records, including failures, rather than selecting only favorable cases.
